\documentclass{article}
\usepackage[margin=1in]{geometry}
\begin{document}

\title{Spaced Practice in Introductory Chemistry}
\maketitle

\begin{abstract}
We test spaced practice quizzes in an introductory chemistry course with 96 students.
The quiz group scored 78.5\% on the final exam, 6.1 points above the comparison group.
\end{abstract}

\section{Introduction}
Spaced practice improves long term retention in many settings \citep{cepeda2006}.
Quizzing also helps learners notice what they do not yet know \citet{roediger2006}.

\section{Method}
We enrolled 96 students ($n = 96$) and assigned sections at random.
Of these, 50\% were in their first year, 30\% were in their second year and 20\% were in later years.

\section{Results}
The quiz group scored 78.5\% on the final exam and the comparison group scored 72.4\%, a gap of 6.1 points.
Table~\ref{tab:main} gives the full results, following the reporting advice in \cite{dunlosky2013}.

\begin{table}[t]
\caption{Final exam scores by group.}
\label{tab:main}
\begin{tabular}{lc}
\toprule
Group & Score (\%) \\
\midrule
Quiz & 78.5 \\
Comparison & 72.4 \\
\bottomrule
\end{tabular}
\end{table}

\bibliographystyle{plain}
\bibliography{refs}
\end{document}
